% article class because we want to fully customize the page and not use a cv template
\documentclass[a4paper,11pt]{article}

%----------------------------------------------------------------------------------------
%FONT
%----------------------------------------------------------------------------------------

% % fontspec allows you to use TTF/OTF fonts directly
% \usepackage{fontspec}
% \defaultfontfeatures{Ligatures=TeX}

% % modified for ShareLaTeX use
% \setmainfont[
% SmallCapsFont = Fontin-SmallCaps.otf,
% BoldFont = Fontin-Bold.otf,
% ItalicFont = Fontin-Italic.otf
% ]
% {Fontin.otf}

%----------------------------------------------------------------------------------------
%PACKAGES
%----------------------------------------------------------------------------------------
\usepackage{url}
\usepackage{parskip} 
\usepackage[a4paper,left=1cm,right=1cm,top=1.0cm,bottom=1.0cm]{geometry}
%other packages for formatting
\RequirePackage{color}
\RequirePackage{graphicx}
\usepackage[usenames,dvipsnames]{xcolor}
\usepackage[scale=0.9]{geometry}

%tabularx environment
\usepackage{tabularx}

%for lists within experience section
\usepackage{enumitem}
\setlist[itemize]{label=\textcolor{black}{\textbullet}}
\setlist[itemize]{font=\color{black}}
\renewcommand{\labelitemi}{\textcolor{black}{\textbullet}}

% centered version of 'X' col. type
\newcolumntype{C}{>{\centering\arraybackslash}X} 

%to prevent spillover of tabular into next pages
\usepackage{supertabular}
\usepackage{tabularx}
\newlength{\fullcollw}
\setlength{\fullcollw}{0.47\textwidth}

%custom \section
\usepackage{titlesec}
\usepackage{multicol}
\usepackage{multirow}

%CV Sections inspired by: 
%http://stefano.italians.nl/archives/26
\titleformat{\section}{\Large\scshape\raggedright}{}{0em}{}[\titlerule]
\titlespacing{\section}{0pt}{6pt}{6pt}

%for publications
\usepackage[style=authoryear,sorting=ynt, maxbibnames=2]{biblatex}

%Setup hyperref package (Active links disabled via draft=true)
\usepackage[unicode, draft=true]{hyperref}
% \definecolor{linkcolour}{rgb}{0,0.2,0.6}
% \hypersetup{colorlinks,breaklinks,urlcolor=linkcolour,linkcolor=linkcolour}
\addbibresource{citations.bib}
\setlength\bibitemsep{1em}

%for social icons
\usepackage{fontawesome5}

%debug page outer frames
%\usepackage{showframe}


% job listing environments
\newenvironment{jobshort}[2]
    {
    \begin{tabularx}{\linewidth}{@{}l X r@{}}
    \textbf{#1} & \hfill &  #2 \\[3.75pt]
    \end{tabularx}
    }
    {
    }

\newenvironment{joblong}[2]
    {
    \begin{tabularx}{\linewidth}{@{}l X r@{}}
    \textbf{#1} & \hfill &  #2 \\[3.75pt]
    \end{tabularx}
    \begin{minipage}[t]{\linewidth}
    \begin{itemize}[nosep,after=\strut, leftmargin=1em, itemsep=3pt,label=--]
    }
    {
    \end{itemize}
    \end{minipage}    
    }


%----------------------------------------------------------------------------------------
%BEGIN DOCUMENT
%----------------------------------------------------------------------------------------
\begin{document}

% non-numbered pages
\pagestyle{empty} 

%----------------------------------------------------------------------------------------
%TITLE
%----------------------------------------------------------------------------------------

\begin{tabularx}{\linewidth}{@{}C@{}}
\Huge{Muk Chunpongtong} \\[7.5pt]

\raisebox{-0.05\height}\ github.com/muxite \ $|$ \ 
\raisebox{-0.05\height}\ muk.chunpongtong@gmail.com \ $|$ \ 
\raisebox{-0.05\height}\ +1\ 825\ 558\ 0137 \\

\end{tabularx}


%----------------------------------------------------------------------------------------
%EDUCATION
%----------------------------------------------------------------------------------------

\section{Education}

\begin{tabularx}{\linewidth}{@{}l X r@{}}
\textbf{University of British Columbia (2nd Year)} & \hfill & Sept 2024 - May 2029 \\[3.75pt]
\end{tabularx}
\begin{minipage}[t]{\linewidth}
Bachelor of Applied Science, Computer Engineering \hfill \textbf{Overall average: 84.2\%}
\end{minipage}

\vspace{-1mm}

%----------------------------------------------------------------------------------------
%TECHNICAL SKILLS
%----------------------------------------------------------------------------------------
\section{Technical Skills}

\begin{tabularx}{\linewidth}{@{}>{\raggedright\arraybackslash}p{3cm} X@{}}
Languages & Python, C/C++, Java, SQL, MATLAB, SystemVerilog, CUDA, RISC-V Assembly \\
Libraries & OpenCV, Tensorflow, Selenium, SFML, Asyncio, PyTest, FastAPI, ReactJS \\
Tools & Docker, AWS ECS, CMake, Bazel, Kubernetes, RabbitMQ, FPGA, ESP32, Quartus \\
General Skills & Git, Vector Databases, Oscilloscope, Linux \& Windows CLI, SSH \\
\end{tabularx}


%----------------------------------------------------------------------------------------
%WORK EXPERIENCE
%----------------------------------------------------------------------------------------

\section{Work Experience}

\begin{joblong}{Teaching Assistant (APSC 160), UBC CS Department}{Oct 2025 – Jan 2026}
\item Led weekly office hours conducting live code reviews and debugging sessions to teach problem-solving strategies.
\item Co-hosted lab sessions for 30 students, delivering C programming mini-lectures and facilitating live Q\&A.
\item Evaluated 300+ C assignments by manually testing code in Arduino simulator and providing targeted feedback.
\end{joblong}

\vspace{-7mm}


%----------------------------------------------------------------------------------------
%TECHNICAL PROJECTS
%----------------------------------------------------------------------------------------

\section{Technical Projects}

\begin{joblong}{WebRAG, Autonomous Research Assistant (Full-Stack) \textit{(Personal Project)}}{Jul 2025 -- Present}
\item Developed a scalable agentic research and reasoning assistant that decomposes queries into parallel web actions and returns structured, source-cited answers with a full website frontend.
\item Designed a Graph-of-Thoughts planner with score-guided pruning, improving task accuracy by +26.8\% while reducing token usage by -29.3\% versus a sequential baseline.
\item Built a benchmark runner + GDB-inspired debugger to drive experiment-based iteration on up to 34 custom scenarios, with hybrid scoring (links, key terms, LLM rubric).
\item Added end-to-end observability across connectors and DAG nodes, logging all inputs/outputs to trace each step.
\item Deployed on AWS: built a backend with Redis, a vector DB, ECS workers, and a custom Lambda-based autoscaler to handle concurrent requests and persistent task state.
\end{joblong}




\begin{joblong}{WaveSim, C++/CUDA Wave Simulation \textit{(Personal Project)}}{May 2025 - Present}
\item Built a real-time 2D wave physics engine in C++ to simulate fluid dynamics and wave propagation for interactive visualization, using SFML for rendering and CMake for cross-platform builds.
\item Accelerated the simulation with GPU parallel processing via CUDA, implementing the discrete wave equation and achieving 39× speedup over the CPU implementation.
\end{joblong}

\begin{joblong}{AutoSpotCV, Face Recognition Pipeline \textit{(Personal Project)}}{Jun 2024 -- Aug 2024}
\item Built a Python face-recognition pipeline in OpenCV with automated data collection, data augmentation, and real-time inference on a custom dataset.
  \item Implemented LBPH and a TensorFlow classifier; analyzed per-subject sensitivity scores (0.20--0.85) to identify error-prone identities and guide data collection and augmentation.
  \item Automated collection with Selenium, capturing 1000+ labeled images and cutting manual prep by 91\%.
\end{joblong}



%----------------------------------------------------------------------------------------
%ENGINEERING OR DESIGN STUDENT TEAMS
%----------------------------------------------------------------------------------------


\section{Extracurricular Student Teams}

\begin{joblong}{Software Developer, UBC ThunderBots RoboCup Team \textit{(Design Team)}}{Sept 2024 -- Present}
\item Build autonomous soccer-playing robots for international RoboCup SSL, contributed to gameplay AI and robot control in the team simulation + on-robot stack.
\item Refactored gameplay from coroutines to FSMs in C++ (Bazel), improving maintainability and deploying smarter ball-targeting in simulation and live trials.
\item Fixed 3 simulated Python tests, raising branch coverage by 10\% across the decision-making pipeline.
\item Collaborated via reviews/tickets/PRs with a 20-person team to close coverage gaps.
\end{joblong}

%\begin{joblong}{President \& Cofounder, Queen Elizabeth High School Engineering Club}{Sept 2022 -- Jun 2024}
%\item Founded and led a 40-student engineering club, planning and running hands-on projects and workshops.
%\item Secured hundreds of dollars in school funding by pitching project proposals to faculty and the principal; delivered planters and composters with club members, expanding the community garden by 15\%.
%\end{joblong}


\end{document}
